\documentclass[11pt, letterpaper]{article}

\usepackage[margin=1in]{geometry}
\usepackage{amsmath}
\usepackage{amssymb}
\usepackage{booktabs}
\usepackage{array}
\usepackage{hyperref}
\usepackage{xcolor}
\usepackage{listings}
\usepackage{caption}
\usepackage{parskip}
\usepackage{titlesec}
\usepackage{microtype}

\hypersetup{colorlinks=true, linkcolor=blue, urlcolor=blue, citecolor=blue}

\definecolor{codegray}{rgb}{0.95,0.95,0.95}
\lstset{
  backgroundcolor=\color{codegray},
  basicstyle=\ttfamily\small,
  breaklines=true,
  frame=single,
  rulecolor=\color{gray!40}
}

\title{\textbf{Diffusion Accelerator: An Address-Centric RTL Architecture\\
for Accelerating Stable Diffusion Inference}\\[0.4em]
\large Implementation and Evaluation Report}
\author{}
\date{\today}

\begin{document}
\maketitle
\tableofcontents
\newpage

% ─────────────────────────────────────────────────────────────────────────────
\section{Introduction}
% ─────────────────────────────────────────────────────────────────────────────

Stable Diffusion models rely on a U-Net backbone that alternates between two
computationally distinct operation types: spatial convolution layers in the
encoder and decoder, and self-attention layers (implemented as matrix
multiplications) in the bottleneck~\cite{sdacc2025}. When executing this
workload on a general-purpose GPU, two hardware-level bottlenecks dominate
cycle counts:

\begin{enumerate}
  \item \textbf{Im2col buffer overhead~\cite{sdacc2025}.} GPUs are
        fundamentally GEMM (General Matrix Multiply) engines. Convolution's
        irregular patch-access pattern must be physically reformatted into a
        dense matrix before any computation can begin. This intermediate
        buffer materialisation is pure overhead — the original data already
        resides in memory and must be re-read.

  \item \textbf{Heterogeneous operation switching~\cite{sdacc2025}.} The
        transition between convolutional layers (\texttt{cuDNN}) and
        attention/matmul layers (\texttt{cuBLAS}) requires a full kernel
        context switch on a GPU, incurring hundreds of wasted cycles per
        transition.
\end{enumerate}

This project implements a synthesisable RTL accelerator in SystemVerilog that
addresses both bottlenecks. The design is verified cycle-accurately using
Verilator~5.048, with all 256 INT32 output values confirmed correct against a
C++ golden reference model.

% ─────────────────────────────────────────────────────────────────────────────
\section{Architecture Overview}
% ─────────────────────────────────────────────────────────────────────────────

The accelerator consists of four primary RTL modules that together replace the
im2col+GEMM pipeline with a direct, address-centric dataflow.

\subsection{Systolic Array (\texttt{systolic\_array\_4x4.sv})}

A $4 \times 4$ array of processing elements (PEs), each performing an INT8
multiply-accumulate (MAC) operation and accumulating into a 32-bit register.
Data flows in a skewed (systolic) pattern: input activations are fed along
rows and weights along columns, with each PE passing its input to its
right-neighbour and down-neighbour on the following cycle. The array is
general-purpose and serves as the shared compute substrate for both
convolution and matrix multiply.

\subsection{Convolution Address Generator (\texttt{conv\_addr\_gen.sv})}

The core contribution of the design. Rather than pre-computing an im2col
buffer, this module computes the required input SRAM address \emph{in the
same cycle the data is needed}, using purely combinational logic:

\begin{equation}
  i_h = o_h \cdot s + k_r - p, \qquad
  i_w = o_w \cdot s + k_c - p
  \label{eq:coord}
\end{equation}

\begin{equation}
  \text{addr} = (i_c \cdot H_{\text{in}} + i_h) \cdot W_{\text{in}} + i_w
  \label{eq:addr}
\end{equation}

where $o_h, o_w$ are the output pixel coordinates, $k_r, k_c$ the kernel
offsets, $i_c$ the input channel index, $s$ the stride, and $p$ the padding.
The module also asserts \texttt{addr\_valid = 0} for out-of-bounds accesses
(i.e., zero-padding regions), gating the SRAM read enable so that no memory
access is wasted on padding values. The intermediate buffer is never
allocated.

\subsection{Convolution Controller (\texttt{conv\_controller.sv})}

An FSM-driven controller that orchestrates the systolic array for 2D
convolution. The state machine progresses through five phases:

\begin{center}
\texttt{IDLE} $\to$ \texttt{PRELOAD} $\to$ \texttt{FEED} $\to$
\texttt{DRAIN} $\to$ \texttt{OUTPUT}
\end{center}

During \texttt{PRELOAD}, weight values are loaded into the PE array. During
\texttt{FEED}, the address generator streams input activations directly into
the array one pixel at a time. The \texttt{DRAIN} phase flushes the pipeline
after the last activation is presented.

\subsection{Matmul Controller (\texttt{matmul\_controller.sv})}

A second FSM controller for general matrix multiplication, also targeting the
shared systolic array. Address generation here is sequential and trivial.
This controller enables the same hardware to service both convolution and
attention layers in the U-Net, with no structural modifications to the array
itself.

\subsection{Top-Level Accelerator (\texttt{top\_accel.sv})}

A top-level module that exposes a \texttt{mode} register. When
\texttt{mode = CONV}, the conv controller drives the systolic array; when
\texttt{mode = MATMUL}, the matmul controller drives it. The switch requires
a single register write and takes effect within 5 cycles, compared to the
hundreds of cycles required for a GPU kernel context switch.

% ─────────────────────────────────────────────────────────────────────────────
\section{GPU Baseline Model}
% ─────────────────────────────────────────────────────────────────────────────

To produce a meaningful cycle-count comparison, a cycle-accurate model of the
traditional GPU convolution pipeline was constructed. For an
$8 \times 8 \times 4$ input with a $3 \times 3$ kernel and 4 output channels,
the GPU baseline is decomposed as follows:

\begin{table}[h]
\centering
\caption{GPU baseline cycle breakdown for a single convolution layer}
\begin{tabular}{lrr}
\toprule
\textbf{Phase} & \textbf{Cycles} & \textbf{Description} \\
\midrule
Im2col buffer fill  & 2,304  & One cycle per byte written to the intermediate buffer \\
GEMM computation    & 9,217  & Matrix multiply over the 2,304-entry buffer \\
Output write        & 256    & One cycle per output value written to SRAM \\
\midrule
\textbf{Total}      & \textbf{11,777} & \\
\bottomrule
\end{tabular}
\end{table}

The im2col buffer has $8 \times 8 \times (4 \times 3 \times 3) = 2{,}304$
entries: one per output pixel (64) times one per kernel unrolling position
(36). Each weight must be re-read once per output pixel, resulting in
$36 \times 4 \times 64 = 9{,}216$ weight re-reads during GEMM — a $32\times$
redundancy over the accelerator's single-pass weight loading.

% ─────────────────────────────────────────────────────────────────────────────
\section{Verification Methodology}
% ─────────────────────────────────────────────────────────────────────────────

Functional correctness was established using Verilator~5.048, an
open-source RTL simulator that compiles SystemVerilog to C++ and executes it
as a native binary. The testbench proceeds as follows:

\begin{enumerate}
  \item A C++ golden reference model computes the expected INT32 output for
        a $3 \times 3$ convolution on an $8 \times 8 \times 4$ input with
        randomised weights and activations.
  \item The same inputs are loaded into the accelerator's simulated SRAM.
  \item The RTL simulation runs to completion and all 256 output values are
        read from the output SRAM.
  \item Each output is compared to the golden model: all 256/256 values
        match exactly.
\end{enumerate}

Cycle counts are extracted from the simulation trace and exported to
\texttt{trace\_data.js} for replay in the browser-based visualiser.

% ─────────────────────────────────────────────────────────────────────────────
\section{Results}
% ─────────────────────────────────────────────────────────────────────────────

\subsection{Single-Layer Convolution Speedup}

\begin{table}[h]
\centering
\caption{Single-layer conv2d comparison: GPU (im2col) vs.\ SD-Acc}
\begin{tabular}{lrrr}
\toprule
\textbf{Metric} & \textbf{GPU (im2col)} & \textbf{SD-Acc} & \textbf{Reduction} \\
\midrule
Total cycles          & 11,777  & 3,298   & $3.57\times$ \\
Intermediate buffer   & 2,304 B & 0 B     & $100\%$ \\
Weight reads          & 9,216   & 288     & $32\times$ \\
Input SRAM reads      & 1,936   & $\sim$1,620 & $\sim$16\% \\
Outputs written       & 256     & 256     & — \\
\bottomrule
\end{tabular}
\end{table}

The $3.57\times$ cycle speedup derives entirely from eliminating the im2col
buffer materialisation. The address generator produces the correct SRAM
address for every activation in $O(1)$ combinational logic, so no pre-copy
phase is required.

\subsection{Multi-Layer Heterogeneous Speedup}

To demonstrate the second bottleneck, a three-layer U-Net slice was modelled:
Conv $\to$ Self-Attention $\to$ Conv. The GPU model includes a kernel-switch
stall of 800 cycles at each layer boundary (cuDNN $\to$ cuBLAS and back).

\begin{table}[h]
\centering
\caption{3-layer U-Net slice: GPU vs.\ SD-Acc total cycle counts}
\begin{tabular}{lrr}
\toprule
\textbf{Segment} & \textbf{GPU cycles} & \textbf{SD-Acc cycles} \\
\midrule
Conv Layer 1              & 11,777  & 3,298  \\
Context switch (conv$\to$matmul) & 800     & 5      \\
Attention Matmul          & 4,096   & 1,200  \\
Context switch (matmul$\to$conv) & 800     & 5      \\
Conv Layer 3              & 11,777  & 3,298  \\
\midrule
\textbf{Total}            & \textbf{29,250} & \textbf{7,806} \\
\midrule
Speedup                   & \multicolumn{2}{c}{$3.75\times$} \\
Context-switch overhead   & 1,600 cyc (5.5\%) & 10 cyc (0.13\%) \\
\bottomrule
\end{tabular}
\end{table}

The SD-Acc mode register flip requires 5 cycles versus 800 cycles for the GPU
kernel context switch — a $160\times$ reduction in switching overhead. Over a
full U-Net inference pass with tens of layer boundaries, this compounds
significantly.

\subsection{Summary of Key Findings}

\begin{itemize}
  \item Eliminating the im2col buffer yields a $\mathbf{3.57\times}$ speedup
        on a single convolution layer, verified directly from RTL simulation.
  \item Weight re-read count is reduced by $\mathbf{32\times}$ because weights
        are loaded once during PRELOAD and never re-read, whereas GEMM re-reads
        each weight for every output pixel.
  \item Mode switching between convolution and matrix multiply costs
        $\mathbf{5}$ cycles on the accelerator versus $\sim$800 cycles on a
        GPU, a $\mathbf{160\times}$ improvement.
  \item A single $4\times4$ systolic array services both operation types with
        no structural change, reducing area versus a dual-datapath design.
  \item Combined single-layer and switching savings yield a
        $\mathbf{3.75\times}$ end-to-end speedup over a 3-layer U-Net slice.
\end{itemize}

% ─────────────────────────────────────────────────────────────────────────────
\section{Module Summary}
% ─────────────────────────────────────────────────────────────────────────────

\begin{table}[h]
\centering
\caption{RTL module inventory}
\begin{tabular}{lll}
\toprule
\textbf{File} & \textbf{Role} & \textbf{Key property} \\
\midrule
\texttt{systolic\_array\_4x4.sv} & $4\times4$ INT8 systolic array & Shared by conv and matmul \\
\texttt{conv\_addr\_gen.sv}      & Combinational address generator & Zero-cycle im2col replacement \\
\texttt{conv\_controller.sv}     & Conv2d FSM controller & 5-state pipeline \\
\texttt{matmul\_controller.sv}   & Matmul FSM controller & Sequential address generation \\
\texttt{top\_accel.sv}           & Unified top-level & Mode register mux \\
\bottomrule
\end{tabular}
\end{table}

% ─────────────────────────────────────────────────────────────────────────────
\section{Toolchain}
% ─────────────────────────────────────────────────────────────────────────────

\begin{itemize}
  \item \textbf{RTL language:} SystemVerilog (synthesisable subset)
  \item \textbf{Simulator:} Verilator 5.048 (cycle-accurate, compiles to C++)
  \item \textbf{Golden reference:} C++ model (INT32 fixed-point arithmetic)
  \item \textbf{Trace extraction:} Python (JSON generation for visualiser)
  \item \textbf{Visualiser:} HTML5 Canvas + JavaScript (cycle-accurate replay)
\end{itemize}

% ─────────────────────────────────────────────────────────────────────────────
\section{Limitations and Future Work}
% ─────────────────────────────────────────────────────────────────────────────

The current prototype targets an $8 \times 8 \times 4$ feature map with a
$4\times4$ systolic array. Real Stable Diffusion layers operate at
$64\times64$ with 320+ channels, requiring tiling logic and a larger array
(e.g., $16\times16$ or $32\times32$) to handle full-resolution inference.

The phase-aware sampling optimisation described in~\cite{sdacc2025} — which
identifies and skips redundant diffusion timesteps at the algorithmic level —
was not implemented in this project. Adding this would compound the hardware
speedup with an algorithmic one and represents the most impactful near-term
extension.

Longer-term, a full chip specification targeting a 28\,nm process node, with
on-chip SRAM banks sized for a complete U-Net encoder block and a DDR memory
controller for weight streaming, would be required to evaluate the design in a
realistic system context.

% ─────────────────────────────────────────────────────────────────────────────
\bibliographystyle{plain}
\begin{thebibliography}{9}

\bibitem{sdacc2025}
Anonymous.
\newblock {SD-Acc: Accelerating Stable Diffusion through Phase-aware Sampling
and Hardware Co-Optimizations}.
\newblock \textit{arXiv preprint arXiv:2507.01309}, 2025.
\newblock \url{https://arxiv.org/abs/2507.01309}

\end{thebibliography}

\end{document}
